% !TEX root = main.tex
\documentclass[main.tex]{subfiles}
\begin{document}

\section{Results}
\label{sec:results}

\subsection{Inverse Tessellations}
\label{sec:results-inverse}

% TODO: compare our joint and softened formulations (Sec. 4) against the inverse/fitting examples of Bourne et al.
% Test cases from the papers (in bib/):
% - [Bourne2025Inverting, Sec. 5.1, Figs. 1-2]: recover a Laguerre diagram from the areas and centroids of its cells (2D, random generators).
% - [Bourne2025Inverting, Sec. 5.2, Figs. 3-7]: fit to synthetic data perturbed with small (0.001) and larger (0.05) noise; minimum-distance constraint active.
% - [Bourne2025Inverting, Sec. 5.3, Figs. 8-9]: fit to an EBSD image of steel.
% - [Bourne2026ConvexOrder, Sec. 6.3.1, Figs. 8-14]: reconstruction with N = 20 and N = 728 on [0,1]^2, t = 1250; subgradient vs Frank-Wolfe convergence.
% - [Bourne2026ConvexOrder, Sec. 6.3.2, Figs. 15-16]: EBSD image of single-phase steel, N = 243 grains, domain [0, 252.25]^2, t = 10.
% Metrics to report: centroid residual, volume error, iterations / volume solves, wall-clock time, and failures (merged generators).

\subsection{Incompressible Flows}
\label{sec:results-flows}

% TODO: compare the constrained (Sec. 5.2) and penalty-based (Sec. 5.3) integrators against Power Particles [deGoes2015PowerParticles] and the scheme of [Gallouet2018Lagrangian].
% Metrics to report: energy drift over time, centroid--particle gap, volume error, divergence, Newton iterations per step, maximum stable time step, and sensitivity to alpha and epsilon.

\end{document}
